\documentclass[12pt]{article}

\usepackage[utf8]{inputenc}
\usepackage[T2A]{fontenc}
\usepackage[russian]{babel}

\usepackage{amsmath,amssymb,amsthm}
\usepackage{graphicx}
\usepackage{mathtools}
\usepackage{listings}

\usepackage[backend=biber, style=gost-numeric, sorting=none]{biblatex}
\addbibresource{references.bib}

\newtheorem{definition}{Definition}
\newtheorem{theorem}{Theorem}
\newtheorem{example}{Example}[section]
\newtheorem{remark}{Remark}

\begin{document}

\renewcommand{\baselinestretch}{1.5}
\renewcommand{\contentsname}{Содержание}

\tableofcontents
\newpage

\section{Модель GTN и предварительные сведения}
Главной темой предыдущей лекции была теорема Чена и Чена об универсальной аппроксимации нелинейных операторов в пространствах непрерывных функций.
Соответствующая нейронная сеть имеет архитектуры вида \(N \circ L\) (выход линейного оператора сэмплирования/проецирования передаётся в нелинейную часть) и была построена с использованием функции Таубера--Винера в качестве активационной, что и позволило добиться успеха.
Однако этот результат носит больше абстрактный характер, Мхаскар и Хахм (Mhaskar \& Hahm) сделали переход между \(N\) и \(L\) конструктивным с почти оптимальным количеством параметров.

В основном будем работать на \(s\)-мерном кубе \([-1, 1]^s\), \(s \in \mathbb{N}\) с Гёлдеровской нормой \(\|\cdot\|_p\).
Для \(\lambda \ge 0\) через \(\Pi_{\lambda,s}\) обозначим множество многочленов от \(s\) переменных с покоординатной степенью не больше \(\lambda\).
Часто компактные подмножества \(L^p\) задаются равномерными ограничениями на \emph{наилучшее приближение} многочленами
\[E_{\lambda,p,s}(f) \coloneq \inf_{P \in \Pi_{\lambda,s}} \|f - P\|_p\]
или условиями на коэффициенты Соболева/Лежандра.

Напомним, что \textbf{модулем непрерывности} называют функцию \(\Omega(h) \to 0\) при \(h \to 0\), убывающую и субаддитивную: \(\Omega(h_1 + h_2) \le \Omega(h_1) + \Omega(h_2)\). (Например, \(\Omega(h) = h^\alpha\) с \(0 < \alpha \le 1\).)
По нему можем ввести класс \textbf{равномерно непрерывных функционалов}
\[\mathcal{F}_{\Omega,p,s} \coloneq \left\{ F\colon L^p \to \mathbb{R} : |F(f) - F(g)| \le \Omega(\|f - g\|_p) \ \forall f,g \right\}.\]
Этот инструментарий можно рассматривать как операторный аналог Гёльдеровской или Липшицевой регулярности, используемой для УА функций.
\begin{definition}[Обобщённая сеть переноса, GTN]
  Фиксируем натуральные \(N\) и \(1 \le d \le D\), активационную функцию \(\varphi \colon \mathbb{R}^d \to \mathbb{R}\).
  \textbf{Обобщённой сетью переноса (generalized translation network)} с \(N\) блоками вычисляет
  \[x \in \mathbb{R}^{D} \longmapsto \sum_{k=1}^{N} a_k \varphi(A_k x + b_k),\]
  где \(A_k \in \mathbb{R}^{d \times D}\), \(b_k \in \mathbb{R}^d\), \(a_k \in \mathbb{R}\).
  Множество всех таких сетей обозначим как \(\Pi_{\varphi;N,D}\).
\end{definition}
\textit{Частные случаи:} при \(d = 1\) это многослойный перцептрон с одним скрытым слоем, а при \(d = D\) и радиальной \(\varphi\) получаем RBF сеть.

На \([-1, 1]\) с единичным весом определены ортогональные \textbf{многочлены Лежандра}, формула Родрига для них
\[P_n(t) = \frac{1}{2^n n!} \frac{d^n}{dt^n} \left[(t^2 - 1)^n\right] \qquad (n = 0,1,2,\dots).\]
Но нам будет полезнее нормализованная версия
\[\mathrm{IP}_n(t) = \sqrt{\frac{2n+1}{2}} P_n(t), \quad \int_{-1}^1 \mathrm{IP}_n(t) \mathrm{IP}_m(t) \, dt = \delta_{nm}.\]
\textbf{Основные свойства:}
\begin{itemize}
  \item Рекуррентная формула: \((n+1)P_{n+1}(t) = (2n+1)t P_n(t) - n P_{n-1}(t)\), для \(\mathrm{IP}_n\):
        \[\sqrt{n + \tfrac{1}{2}} \, t \, \mathrm{IP}_n(t) = \alpha_{n+1} \mathrm{IP}_{n+1}(t) + \alpha_n \mathrm{IP}_{n-1}(t), \quad \alpha_n = \frac{n}{\sqrt{(2n-1)(2n+1)}}.\]
  \item Чётность: \(\mathrm{IP}_n(-t) = (-1)^n \mathrm{IP}_n(t)\).
  \item Значения в концах: \(P_n(\pm 1) = (\pm 1)^n\), значит \(\mathrm{IP}_n(\pm 1) = \sqrt{\frac{2n+1}{2}} (\pm 1)^n\).
  \item Полнота в \(L^2([-1,1])\): конечные линейные комбинации \(\{\mathrm{IP}_n(t)\}\) образуют всюду плотное множество.
\end{itemize}
Для обобщения на многомерный куб, перемножим многочлены согласно мультииндексу \(k = (k_1,\dots,k_D) \in \mathbb{Z}_+^D\), \(x = (x_1,\dots,x_D) \in [-1,1]^D\) (тензорное произведение):
\[\mathrm{IP}_k(x) \coloneq \prod_{j=1}^D \mathrm{IP}_{k_j}(x_j).\]
Тогда, согласно теореме Фубини, \(\{\mathrm{IP}_k\}_{k \in \mathbb{Z}_+^D}\) образуют ортонормированный базис в \(L^2([-1,1]^D)\), а значит применима общая теория рядов Фурье.
Для \(g \in L^1([-1,1]^D)\) (или \(L^2\)) определим
\[a_k(g) \coloneq \int_{[-1,1]^D} g(t) \mathrm{IP}_k(t) \, dt, \qquad g \sim \sum_{k \in \mathbb{Z}_+^D} a_k(g) \mathrm{IP}_k\]
(со сходимостью в \(L^2\), кроме того верно равенство Парсеваля \(\|g\|_2^2 = \sum_k |a_k(g)|^2\) при \(g \in L^2\)).
Выберем конечное \(\Pi \subset \mathbb{Z}_+^D\), определяющее какие элементы базиса сохраняться в усечённом разложении, оператор проецирования на \(\mathrm{span}\{\mathrm{IP}_k : k \in \Pi\}\):
\[(P_\Pi g)(x) = \sum_{k \in \Pi} a_k(g) \mathrm{IP}_k(x).\]
\begin{example}
  Некоторые классические варианты \(\Pi\) и их мощности.
  \begin{table}[!h]
    \centering
    \begin{tabular}{|r|r|l|}
      \hline
      Обозначение            & Условие на \(k\)                         & Мощность                                     \\
      \hline
      \(\Pi_s^{\text{tot}}\) & \(k_1 + \cdots + k_D \le s\)             & \(C_{D + s}^s\)                              \\
      \(\Pi_s^{\max}\)       & \(\max_{1 \le i \le D}{k_i} \le s\)      & \((s + 1)^D\)                                \\
      \(\Pi_3^{\text{HC}}\)  & \(\prod_{j = 1}^{D}(k_j + 1) \le s + 1\) & \(\sim\frac{(s+1)[\ln(s+1)]^{D-1}}{(D-1)!}\) \\
      \hline
    \end{tabular}
  \end{table}
\end{example}

Для узлов \(X = \{x_j\}_{j=1}^m \subset [-1,1]^D\), определим \textbf{оператор сэмплирования} \(L\colon \{x_j\}_{j=1}^m \subset [-1,1]^D\), \(L(g) = (g(x_1), \dots, g(x_m))\).
По нему можем построить \textbf{коэффициенты приближения (МНК/квадратуры)}.
Пусть \(V_{j,k} = \mathrm{IP}_k(x_j)\), \(W > 0\) --- матрица весов, тогда имеем линейную оценка
\[\hat{a}(g) \coloneqq (V^T W V)^{-1} V^T W L(g) \in \mathbb{R}^{|\Pi|}, \qquad \tilde{a}_k(g) \approx a_k(g).\]
Этот вектор будем использовать в \textbf{полиномиальном заменителе}: для вектора \(a \in \mathbb{R}^{|\Pi|}\) положим
\[I_\Pi(a)(x) = \sum_{k \in \mathbb{N}} a_k \mathrm{IP}_k(x).\]

\begin{remark}[О выборе \(\Pi\)]
  Практические советы по выбору набора мультииндексов \(\Pi\), чтобы не потерять в точности и времени вычислений:
  \begin{itemize}
    \item Больше \(\Pi \Rightarrow\) меньше смещение из-за усечения, но больше \(|\Pi|\), что может приводить к плохой обусловленности.
    \item Для гладких/аналитических \(g\) часто \(\Pi_s^{\text{tot}}\) оптимальный выбор, для больших \(D\) \(\Pi_s^{\text{HC}}\) уменьшает взрыв размерности.
  \end{itemize}
  \underline{Инструкция по выбору (в стиле Мхаскара--Хахма).}
  \begin{enumerate}
    \item Зафиксировать семейство индексов \(\Pi_s\) и увеличивать \(s\), пока ошибка усечения больше заданного значения.
    \item Для каждого \(\Pi_s\) выбрать \(m \gtrsim c|\Pi_s|\) точек и построить \(I_{\Pi_s}(\hat{a}(g))\).
    \item Следить за ошибкой на тестовой сетке, чтобы остановиться, когда \(\|g - I_{\Pi_s}(\hat{a}(g))\|\) имеет заданную точность.
  \end{enumerate}
\end{remark}

Перейдём теперь непосредственно к \underline{построению операторов \(L\) и \(N\)}.\\
\textbf{Линейная часть:} для \(m \in \mathbb{N}\) положим \(D = (m + 1)^s = |\Pi_m^{\max}|\) и \(L\colon L^p([-1, 1]^s) \to \mathbb{R}^D\) будет возвращать все коэффициенты Фурье--Лежандра не выше \(m\)-го порядка по каждой координате, а \(I_\Pi\colon \mathbb{R}^D \to \Pi_{m, s}\) восстанавливает по ним полином.
Тогда имеют места неравенства
\[\|L_m(f)\| \le C \|f\|_p, \qquad \|f - I_\Pi L_m(f)\|_p \le \varepsilon_m(K) \ (f \in K)\]
с явными константами из неравенств типа Джексона/Маркова.\\
\textbf{Нелинейная часть:} GTN \(\sum_k a_k \varphi(A_k \cdot + b_k)\) на \(\mathbb{R}^D\).

\underline{Откуда берутся оценки?}
Поскольку \(\{\mathrm{IP}_k\}\) ортонормированна в \(L^2\) и равномерно ограничена на \([-1, 1]\), найдётся \(C = C(p,s)\), что
\[\|L_m(f)\|_{\ell^p(\Pi_m^{\max})} \le C \|f\|_{L^p([-1,1]^s)} \quad (1 \leq p \leq \infty).\]
Дальше применим \emph{(прямую) теорему Джексона} на \([-1, 1]^s\), для \(r \in \mathbb{N}\) и показателя \(p\)
\[E_{m,p,s}(f) \le C \omega_r \left( f; \tfrac{1}{m} \right)_{L^p},\]
где \(\omega_r\) --- это \(r\)-ый модуль непрерывности.
Если дополнительно \(f \in W^{r, p}\) (пространство Соболева), то
\[E_{m,p,s}(f) \le C m^{-r} \|f\|_{W^{r,p}}.\]
Также нам поможет \emph{(обратное) неравенство Маркова для полиномов}, для \(p_m \in \Pi_{m, s}\) и мультииндекса \(\alpha\)
\[\|D^\alpha p_m\|_{L^p} \le C m^{|\alpha|} \|p_m\|_{L^p}.\]
Теперь соберём это всё вместе, если \(K \subset L^p\) --- компакт с общем модулем гладкости (равностепенная непрерывность, ограниченность в пространстве Соболева), тогда
\[\sup_{f \in K} \|f - I_{\Pi_m^{\max}} L_m(f)\|_{L^p} \leq \varepsilon_{m,K} \xrightarrow[m \to \infty]{} 0.\]

Введём в рассмотрение целевой функционал \(F\colon K \to \mathbb{R}\), непрерывный на компакте \(K\).
Образ \(U_m \coloneq L_m(K) \subset \mathbb{R}^D\) компактен.
На нём определена функция \(\varphi_m(a) \coloneq F(I_{\Pi_m^{\max}}(a))\).
Применим УАТ на \(U_m\): для всякого \(\varepsilon > 0\) существует обобщённая сеть переноса (GTN)
\[N(v) = \sum_{k = 1}^M \alpha_k\varphi(A_k\cdot v + b_k), v \in \mathbb{R}^D,\]
(с неполиномиальной активационной функцией) такая, что
\[\sup_{a \in U_m}|\varphi_m(a) - N(a)| \le \varepsilon.\]
Применяя неравенство треугольник, получаем для всех \(f \in K\)
\[|F(f) - N(L_m(f))| \le \underbrace{|F(f) - F(I_{\Pi_m^{\max}}L_m(f))|}_{\text{линейное приближение}} + \underbrace{|\varphi_m(L_m(f)) - N(L_m(f))|}_{\text{GTN приближение}}.\]
Если \(F\) равномерно непрерывен на \(K\) и \(\|f - I_{\Pi_m^{\max}}L_m(f)\|_p \le \varepsilon_{m, K}\), то первое слагаемое стремится к нулю при \(m \to \infty\), второе мало по УАТ.
Если же \(K\) ограничено в \(W^{r, p}\), то \(\|f - I_{\Pi_m^{\max}}L_m(f)\|_p \le m^{-r}\), что позволяет найти подходящее \(m\).
Аналитичность \(f\) на эллипсе Бернштейна обеспечивает экспоненциальное затухание проекций Лежандра/тригонометрических: \(E_{m,p,s}(f) \lesssim \rho^{-m}\), \(\rho > 1\), то есть можно брать не слишком большое \(m\), что положительно сказывается на \(D = (m + 1)^s\) и уменьшает сложность сети \(N\).

\section{Почти оптимальная универсальная аппроксимация функционалов на \(L^p\)}
Формализуем предыдущие рассуждения следующей теоремой.
\begin{theorem}[Мхаскар--Хахм, теорема 2.1]
  Пусть \(d \ge 1\) и \(\varphi \in C^\infty\) на открытом шаре в \(\mathbb{R}^d\) и \(D^k \varphi(b) \neq 0\) для всех мультииндексов \(k \in \mathbb{Z}_+^d\) и некоторой точки \(b\).
  Пусть \(s \ge 1\), \(1 \le p \le \infty\), \(K \subset L^p([-1,1]^s)\) --- компакт, \(\Omega\) --- модуль непрерывности.
  Положим \(D = (m+1)^s\), \(L_m\) как ввели ранее.
  Тогда существуют линейные функционалы \(\gamma_k\colon C(K) \to \mathbb{R}\) \((k = 1,\dots,N)\) и матрицы \(A_k \in \mathbb{R}^{d \times D}\) ранга \(d\), c \(N\), \(m\), \(n\) зависящими лишь от \((p,s,\varphi,K,\Omega)\), что для любого \(F \in \mathcal{F}_{\Omega,p,s}\) справедливо
  \[\sup_{f \in K} \left| F(f) - \sum_{k=1}^{N} \gamma_k(F) \varphi(A_k L_m(f) + b) \right| \leq c \left\{ \Omega(\varepsilon_{m,K}) + \Omega(m^\beta / n) \right\},\]
  где \(\beta > 0\) --- константа роста Маркова/Джексона (зависит от \(p\), \(s\)) и \(n\) --- степень многочлена, использованного на \(U_m\).
  Архитектура и \(L_m\) не зависят от \(F\), лишь \underline{внешние} коэффициенты зависят линейно от него.
\end{theorem}
\begin{proof}
  Из теорем Джексона/Маркова на \([-1, 1]^s\) (см. предыдущий пункт) для компактного \(K\) с общим модулем гладкости/непрерывности имеем
  \[\sup_{f \in K} \|f - I_{\Pi_m^{\max}} L_m(f)\|_{L^p} \le \varepsilon_{m,K} \xrightarrow[m \to \infty]{} 0.\]
  Поскольку \(F \in \mathcal{F}_{\Omega,p,s}\), то
  \[\sup_{f \in K} |F(f) - F(I_{\Pi_m^{\max}} L_m(f))| \le \Omega(\varepsilon_{m,K}).\]
  Дальше приблизим \(\mu_m(F, L_m(\cdot)) \coloneq F(I_{\Pi_m^{\max}} L_m(\cdot))\) равномерно на \(K\).

  \(U_m\) является компактом в \(\mathbb{R}^D\).
  Для \(a, d \in U_m\) справедливо
  \[\|I_{\Pi_m^{\max}}(a) - I_{\Pi_m^{\max}}(d)\|_{L^p} \le C_m \|a - d\|_{\ell^p}\]
  с константой \(C_m \lesssim m^\beta\) (обратное неравенство Маркова на \(\Pi_{m,s}\)).
  Благодаря модулю непрерывности \(F\)
  \[|\mu_m(F,a) - \mu_m(F,d)| \le \Omega(C_m \|a - d\|_{\ell^p}) \le \Omega(c m^\beta \|a - d\|_{\ell^p}),\]
  значит \(\mu_m(F,\cdot)\) тоже равномерно непрерывен на \(K\).
  По многомерной теореме Джексона или Ньюмана--Шапиро (приближение многочленами, управляемая модулем непрерывности) найдётся полином \(p_n\) с \(\operatorname*{deg} p_n \le n\), обеспечивающий
  \[\sup_{a \in U_m} |\mu_m(F,a) - p_n(a)| \le c \, \Omega(m^\beta / n).\]
  \emph{Идея доказательства:} покрыть \(U_m\) кубами со стороной \(\frac{1}{n}\), применить наследованный модуль непрерывности \(h \mapsto \Omega(c m^\beta h)\) к \(\mu_m(F,\cdot)\) и многочленам Бернштейна, чтобы получить оценку.

  Используем, что \(\varphi \in C^\infty\) и \(D^k \varphi(b) \neq 0\) для всех \(k\), чтобы записать многомерную формулу Тейлора в точке \(b\)
  \[\varphi(b + z) = \sum_{\alpha \in \mathbb{Z}_+^d} \frac{D^\alpha \varphi(b)}{\alpha!} z^\alpha.\]
  Для любого \(n\) найдётся конечное число матриц \(A_k\) (ранга \(d\)) и скаляров \(c_{r, \alpha}\), что на ограниченных множествах
  \[\left\| a^\alpha - \sum_{r} c_{r,\alpha} \varphi(A_r a + b) \right\| \le \varepsilon \qquad \forall |\alpha| \leq n.\]
  \emph{Причина:} матрица системы (нижнетреугольная) \(\{D^\beta \varphi(b)\}_{|\beta| \le n}\) обратима по степеням, поэтому можем решить её, чтобы приблизить мономы вплоть до степени \(n\) линейными комбинациями \(\varphi(A \cdot + b)\).

  Пусть \(p_n(a) = \sum_{|\alpha| \leq n} c_\alpha a^\alpha\).
  Используя предыдущее соображение, получаем
  \[p_n(a) \approx \sum_{k} \gamma_k(F) \varphi(A_k a + b).\]
  где \(A_k\) и количество слагаемых зависит лишь от \(n\), \(d\) и \(\varphi\), но не \(F\) или \(K\).
  Чтобы получить \(\gamma_k(F)\), выбираются узлы \(\{a^{(r)}\} \subset U_m\) и многочлены Лагранжа \(\{\ell_r\}\), \(p_n(a) = \sum_r \mu_m(F,a^{(r)}) \ell_r(a)\).
  Каждый \(\ell_r\) представим как комбинацию \(\varphi(A \cdot + b)\), тогда
  \[\gamma_k(F) \coloneq \sum_r \mu_m(F,a^{(r)}) \eta_{k,r} = \sum_r F(h_{\Pi_m^{\max}}(a^{(r)})) \eta_{k,r},\]
  где \(\eta_{k,r}\) коэффициенты из представления \(\{\ell_r\}\). Зависимость функционалов от \(F\) линейная.

  Итого для всех \(f \in K\)
  \begin{multline*}
    \left| F(f) - \sum_{k=1}^{N} \gamma_k(F) \varphi(A_k L_m(f) + b) \right| \le |F(f) - F(p_m(f))| + \\
    + |\mu_m(F, L_m(f)) - p_n(L_m(f))| + \left| p_n(L_m(f)) - \sum_{k} \gamma_k(F) \varphi(A_k L_m(f) + b) \right|.
  \end{multline*}
  Первое слагаемое \(\le \Omega(\varepsilon_{m,K})\), второе \(\le c\Omega(m^\beta / n)\), третье можно сделать \(\le \varepsilon\).
  Соберём последнее \(\varepsilon\) в константу, чтобы получить требуемую оценку.
\end{proof}

\section{Нелинейная \textit{N}-ширина и нижние оценки на размер}
Дальше мы попробуем выяснить, почему полученная аппроксимация действительно оптимальна.
Введём ряд определений.
\begin{definition}[Нелинейная \(N\)-ширина]
  Это величина
  \[\Delta_N(\mathcal{F}) \coloneq \inf_{M_N, \pi} \sup_{F \in \mathcal{F}, f \in K} |F(f) - M_N(\pi(F), f)|,\]
  где \(\pi\colon \mathcal{F} \to \mathbb{R}^N\) и \(M_N\colon \mathbb{R}^N \times K \to \mathbb{R}\) непрерывны.
\end{definition}
Тут \(\pi\) выполняет роль кодировщика, сжимающего функционал в набор из \(N\) чисел, а \(M_N\) декодирует эту информацию и предсказывает \(F(f)\).
Ширина \(\Delta_N\) является минимальной худшей ошибкой, достижимой в такой схеме.
\begin{definition}[Эллипсоид смешанной гладкости]
  Для \(\alpha > 0\) это эллипсоид вида
  \[K \coloneq \left\{ f \in L^2 : \sum_{k \in \mathbb{Z}_+^n} \|k\|_\infty^{2\alpha} |a_k(f)|^2 \le 1 \right\}, \qquad a_k(f) = \langle f, \mathrm{IP}_k \rangle\]
\end{definition}
Эвристика в том, что \(|a_k(f)| \lesssim |k|_\infty^{-\alpha}\), поэтому усечение на уровне \(\|k\|_m \sim m\) приводит к ошибке \(\asymp m^{-\alpha}\).
``Частотный блок'' \(B_m \coloneq \{k : m \le \|k\|_\infty \le 2m\}\) имеет мощность \(\asymp m^s\), а значит имеет \(2^{cm^s}\) знаковых шаблонов, остающихся в \(K\).
Для каждого знакового вектора \(\sigma \in \{\pm 1\}^{B_m}\) обозначим
\[f_\sigma(x) \coloneq c_m \sum_{k \in B_m} \sigma_k \frac{1}{\|k\|_\infty^\alpha} \mathrm{IP}_k(x), \quad c_m > 0 \text{ выбраны, что } \sum_k \|k\|_\infty^{2\alpha} |a_k(f_\sigma)|^2 \le 1.\]
Тогда \(f_\sigma \in K\) и, по ортонормированности
\[\|f_\sigma - f_\tau\|_{L^2}^2 \ge c m^{-2\alpha} |\{ k \in B_m : \sigma_k \neq \tau_k \}| \ge c' m^{-2\alpha} \quad (\sigma \neq \tau).\]
Поэтому \(\mathcal{P}_m \coloneq \{ f_\sigma : \sigma \in \{\pm 1\}^{B_m} \} \subset K\) состоит из элементов попарно отстоящих друг от друга на \(\gtrsim m^{-\alpha}\), и \(|\mathcal{P}_m| = 2^{|B_m|} \geq 2^{c m^s}\).

Пусть все коды лежат в ограниченном множестве \(Q \subset \mathbb{R}^N\) (всегда возможно для компактного множества функционалов).
Для \(\varepsilon > 0\) число элементов \(\varepsilon\)-покрытия оценивается сверху
\[\mathcal{N}(Q, \varepsilon) \le \left( \frac{C \operatorname{diam}(Q)}{\varepsilon} \right)^N.\]
Выберем \(\varepsilon\), чтобы \(|\mathcal{P}_m| = 2^{c m^s} > \mathcal{N}(Q, \varepsilon)\), тогда найдутся два разных \(\varepsilon\)-близких кода.
Из-за равномерной непрерывности \(M_N\) на \(Q \times K\) малые помехи в коде приводят к малым изменениям в выходе, чтобы сделать их большими, используем разделяющий функционал.
Подойдёт функционал \emph{Липшица} \(F_\star(h) \coloneq \langle h, \psi \rangle_{L^2}\) с \(\|\psi\|_{L^2} = 1 \Rightarrow |F_\star(h) - F_\star(g)| \le \|h - g\|_{L^2}\), то есть \(F_\star \in \mathcal{F}_{\Omega,2,s}\) с \(\Omega(t) = t\).
Чтобы \(F_\star(h) - F_\star(g)\) было сравнимо с \(\|f_\sigma - f_\tau\|_{L^2}\), можно взять \(\psi = \frac{f_\sigma - f_\tau}{\|f_\sigma - f_\tau\|_{L^2}}\).
Тогда для близких кодов и произвольного непрерывного декодера \(M_N\)
\[|F_\star(f_\sigma) - M_N(\pi(F_\star), f_\tau)| + |F_\star(f_\tau) - M_N(\pi(F_\star), f_\tau)| \ge c \|f_\sigma - f_\tau\|_{L^2} \gtrsim m^{-\alpha},\]
что и будет ошибкой в худшем случае.
Для итогового варианта надо применить общий модуль непрерывности: \(\geq \Omega(c m^{-\alpha})\).

В конечном счёте \(2^{c m^s} \gtrsim \mathcal{N}(Q, \varepsilon) \lesssim (C / \varepsilon)^N\), то есть \(c m^s \gtrsim N \log(C / \varepsilon)\).
Фиксируем \(\varepsilon\) и разрешаем уравнение относительно \(m\):
\[m \asymp (\log N)^{1/s}.\]
Итоговая оценка
\[\Delta_N(\mathcal{F}_{\Omega,2,s}) \ge c \Omega\left( (\log N)^{-\alpha/s} \right).\]
Это в точности заключение теоремы 2.2 Мхаскара--Хахма, оно показывает, что обсуждаемая конструкция GTN \emph{почти оптимальная} по \(N\) (нижняя оценка из теории информации).
Действительно, выберем \(n \asymp m^\beta\), \(m \asymp (\log N)^{1/s}\), тогда верхняя оценка погрешности аппроксимации \(\lesssim \Omega\left( (\log N)^{-\alpha/s} \right)\) (для \(K\) смешанной гладкости \(\alpha\)).

Посмотрим на конструкцию с \textbf{с точки зрения динамических систем} (связь с Ченом и Ченом).
Пусть \(S_\tau\) будет временным сдвигом на сигналах \(u\colon \mathbb{R} \to \mathbb{R}\), \((S_\tau u)(t) = u(t - \tau)\).
Функционалы/операторы, инвариантные к сдвигу: \(F \circ S_\tau = F\) или \(G(S_\tau u) = S_\tau G(u)\).
Линейное отображение \(L_m\) можно выбрать как:
\begin{itemize}
  \item ортогональная проекция (коэффициенты Лежандра) --- устойчивы к естественным потокам в пространстве признаков;
  \item временное окно \(L(u) = (u(t_j))_{j=1}^D\) --- эквивариантно по отношению к \(S_\tau\) (используется в идентификации).
\end{itemize}
И то, и другое соотносится с подходом \(N \circ L\) Чена и Чена и динамическими моделями Нарандры--Партасарати:
\begin{itemize}
  \item \(L\) выделяет конечные ``состояния'' из траекторий;
  \item \(N\) считывает динамику.
\end{itemize}
TW/GTN плотность обеспечивает репрезентативную устойчивость (сдвиги и растяжения против обобщённых переносов) для ridge/RBF конструкций.


\section{Четыре классические динамические модели (Нарендра и Партасарати) для идентификации и управления}

В своей основополагающей работе 1990 года К.С. Нарандра и К. Партасарати предложили четыре базовые нелинейные модели для идентификации и управления дискретными системами <<вход-состояние-выход>>. Для SISO-систем они записываются следующим образом:
\begin{enumerate}
  \def\labelenumi{\Roman{enumi}.}
  \item \(y_p(k+1) = \sum_{i=0}^{n-1} \alpha_i y_p(k-i) + g\big(u(k), u(k-1), \dots, u(k-m+1)\big)\),
  \item \(y_p(k+1) = f\big(y_p(k), y_p(k-1), \dots, y_p(k-n+1)\big) + \sum_{j=0}^{m-1} \beta_j u(k-j)\),
  \item \(y_p(k+1) = f\big(y_p(k), y_p(k-1), \dots, y_p(k-n+1)\big) + g\big(u(k), u(k-1) \dots, u(k-m+1)\big)\),
  \item \(y_p(k+1) = f\big(y_p(k), \dots, y_p(k-n+1), u(k), \dots, u(k-m+1)\big)\).
\end{enumerate}

Неизвестные нелинейности \(f\) и \(g\) аппроксимируются многослойными перцептронами (MLP). Модель I содержит нелинейность только по входу, модель II --- только по выходу, модель III разделяет нелинейности, а модель IV является общей НАРСС (NARX) моделью.

Эти модели мотивированы аналогичными структурами для линейных систем (ARX) и позволяют естественным образом встроить нейронные сети в динамический контекст. В оригинальной статье приведены примеры идентификации для каждой модели (например, система \(y_p(k+1) = 0.3y_p(k) + 0.6y_p(k-1) + f(u(k))\) для модели I \cite{Narendra1990}).

\section{Представление моделей как частных случаев GTN}

Обобщённые трансляционные сети (GTN) представляют собой композицию линейного оператора сэмплирования \(L\) и нелинейного оператора \(N\): \(P \approx N \circ L\). Модели Нарандры–Партасарати естественно вписываются в эту парадигму.

\begin{itemize}
  \item \textbf{Модели I, II, IV (NARX-типа)} непосредственно соответствуют схеме \(N \circ L\). Линейный оператор \(L_D\) формирует вектор-регрессор из запомненных значений входа и выхода:
        \[
          \phi(k) = \big( y(k-1), \dots, y(k-n_y),\; u(k-1), \dots, u(k-n_u) \big) \in \mathbb{R}^{D}.
        \]
        Нейронная сеть \(\Phi\) выступает в роли нелинейного оператора \(N\), отображая \(\phi(k)\) в предсказание \(\hat{y}(k|k-1)\). Таким образом, для модели IV (и её частных случаев I, II) имеем \(\hat{y}(k+1) = \Phi(\phi(k)) = (N \circ L_D)(u, y, k)\).

  \item \textbf{Модель III} может быть представлена как комбинация двух MLP: \(N_f\) для аппроксимации \(f\) и \(N_g\) для аппроксимации \(g\), работающих параллельно на разных частях регрессора.

  \item \textbf{Модель II (State-Space)} также допускает операторную интерпретацию, но в более общем виде. Система
        \[
          x(k+1) = \Psi(x(k), u(k)), \quad \hat{y}(k) = \Theta(x(k)),
        \]
        где \(\Psi\) и \(\Theta\) --- MLP, представляет собой рекуррентную нейронную сеть. Оператор сдвига во времени \(S\) действует как обобщённая трансляция на пространстве состояний \(x\). Линейное отображение, восстанавливающее состояние из прошлых данных (например, через наблюдатель или энкодер \(E(\phi(k_0))\)), можно рассматривать как часть оператора \(L\), а \(\Theta\) --- как оператор \(N\).
\end{itemize}

Все четыре модели являются частными случаями или вариациями общей архитектуры \(L\)--\(N\), что обеспечивает единый подход к их анализу на основе теорем об универсальной аппроксимации операторов.

\section{Идея замыкания контура обратной связи, замечания про устойчивость и штрафы на основе якобиана/ISS}

При использовании обученной модели в контуре управления (например, в схеме косвенного адаптивного управления или MPC) модель переводится в \emph{параллельный (free-run)} режим, где её собственные предсказания \(\hat{y}\) используются для формирования регрессора на следующих шагах. Это замыкает контур обратной связи и может привести к неустойчивости, даже если исходная система устойчива.

Критически важным становится обеспечение \emph{устойчивости} модели. Для моделей в пространстве состояний (модель II) достаточным условием глобальной устойчивости является \emph{сжатие} (contraction):
\[
  \sup_{x, u} \| \partial_x \Psi(x, u, \theta) \| \leq \lambda < 1.
\]
Более общим условием является \emph{устойчивость вход-состояние (ISS)}: существует функция Ляпунова \(V(x)\) и функции сравнения \(\underline{\alpha}, \overline{\alpha}, \alpha, \sigma\) такие, что
\[
  \underline{\alpha}(\|x\|) \leq V(x) \leq \overline{\alpha}(\|x\|), \quad V(\Psi(x, u)) - V(x) \leq -\alpha(\|x\|) + \sigma(\|u\|),
\]
что гарантирует ограниченность состояния при ограниченном входе.

Для принудительного выполнения этих условий в процессе обучения в функционал потерь добавляют регуляризационные штрафы на якобианы:
\[
  \mathcal{R}(\theta) = \lambda_x \sum_k \|\partial_x \Psi(x(k), u(k))\|_F^2 + \lambda_y \sum_k \|\partial_x \Theta(x(k))\|_F^2 + \lambda_{sn} \sum_\ell (\sigma_{\max}(W_\ell) - \tau)_+^2.
\]

Такая регуляризация не только улучшает устойчивость замкнутого контура, но и стабилизирует вычисление градиентов в BPTT, смягчая проблемы исчезающих и взрывающихся градиентов, возникающих при перемножении матриц Якоби \(\prod_j A(j)\).

Для моделей NARX-типа (I, III, IV) в параллельном режиме также возникают длинные цепочки вычислений, требующие усечённого BPTT с горизонтом \(H\) (обычно \(H \in [5, 20]\)).

\section{Обучающие протоколы и связь с операторным подходом Чен--Чена}

Обучение динамических моделей может проводиться в двух основных режимах:

\begin{enumerate}
  \item \textbf{Последовательно-параллельный (Series-Parallel)} или \emph{teacher forcing}: в процессе обучения в регрессор подставляются измеренные значения выхода \(y(k-i)\). Это разрывает рекуррентные связи и превращает задачу в обучение статического отображения, что позволяет использовать стандартное обратное распространение. Соответствующая функция потерь --- \emph{одношаговая}:
        \[
          J_{\text{sp}}(\theta) = \sum_k \| y(k) - \hat{y}(k|k-1; \theta) \|^2.
        \]
        Этот режим обеспечивает устойчивое обучение (нет накопления ошибки), но может приводить к рассогласованию при переходе в параллельный режим \cite{Narendra1990}.

  \item \textbf{Параллельный (Parallel)} или \emph{free-run}: модель работает автономно, используя собственные предсказания. Обучение требует BPTT по нескольким шагам (горизонт \(H\)) и минимизации \emph{многошаговой} ошибки:
        \[
          J_{\text{par},H}(\theta) = \sum_{k} \sum_{h=1}^{H} \| y(k+h) - \hat{y}_{h\text{-step}}(k; \theta) \|^2.
        \]
        Этот режим лучше соответствует реальной эксплуатации модели в контуре управления, но сложнее в обучении из-за проблем с градиентами. Часто используют гибридный подход: начинают с \(J_{\text{sp}}\), затем дообучают на \(J_{\text{par},H}\).
\end{enumerate}

Операторная теорема Чен--Чена об  предоставляет теоретическое обоснование подхода \(L\)--\(N\). Она утверждает, что используя фиксированные ridge-функции \(g(w_i \cdot x + \theta_i)\) и настраивая только линейные веса \(c_i(f)\) на выходе, можно равномерно аппроксимировать любой непрерывный оператор на компактном множестве входных сигналов.

Модели Нарандры–Партасарати реализуют эту идею на практике: оператор \(L\) (формирование регрессора из задержек) выполняет линейное сэмплирование, а MLP \(\Phi\) реализует нелинейное считывание. Благодаря теореме о плотности MLP в пространстве непрерывных функций, \(\Phi\) может сколь угодно точно аппроксимировать требуемую нелинейность \(F^*\).

Таким образом, модели I–IV представляют собой исторически значимый мост между классической теорией идентификации нелинейных систем и современными операторными подходами на основе нейронных сетей. Их структурная простота, ясная интерпретация в рамках GTN и наличие эффективных алгоритмов обучения (статическое и динамическое обратное распространение) делают их фундаментом для многих современных методов нейросетевого моделирования динамических систем.

Первая часть лекции основана на работе \cite{Mhaskar1997}, а вторая на \cite{Narendra1990}.


\printbibliography

\end{document}